\section{Introduction}
\label{sec:intro}

Large language models (LLMs) have achieved remarkable capabilities across
language understanding, generation, and reasoning tasks. However, they
fundamentally struggle with \emph{deterministic numerical computation}
, specifically the reliable execution of mathematical formulas that require exact
arithmetic on multi-digit numbers
\citep{dziri2023faith}. This limitation is
particularly consequential in professional domains such as financial
underwriting, where metrics like the Debt Service Coverage Ratio (DSCR)
or Loan-to-Value (LTV) must be computed with precision sufficient for
legal and regulatory compliance.

Current approaches to equipping LLMs with computational ability fall into
three categories, each with significant drawbacks:

\begin{itemize}
\item \textbf{External tool calling} (ReAct \citep{yao2023react},
Toolformer \citep{schick2023toolformer}): Each computation requires a
separate inference pass and API round-trip. A chained calculation like
computing Annual Debt Service and then DSCR requires 2+ full
generation cycles.
\item \textbf{Chain-of-Thought prompting} \citep{wei2022chain}: The model
performs arithmetic in its generated text. This works for simple
calculations but systematically fails on six-or-more-digit numbers
and multi-step compound formulas \citep{dziri2023faith}.
\item \textbf{Tool tokens} (ToolkenGPT \citep{hao2023toolkengpt}): Tools
are represented as special vocabulary tokens. Argument extraction
relies on prompting rather than learned representations, and
integration occurs at the language model head rather than within the
model's representational layers.
\end{itemize}

We observe that sparse Mixture-of-Experts (MoE) architectures
\citep{shazeer2017outrageously,fedus2022switch,jiang2024mixtral} provide
a natural mechanism for integrating heterogeneous computation: the expert
routing mechanism already selects specialized processing modules on a
per-token basis. If some of these experts were deterministic functions
instead of neural feed-forward networks, the router could learn
\emph{when} to invoke exact computation and the model could perform it
\emph{within a single forward pass}.

We introduce \textbf{Hybrid MoE}, an architecture that extends the
expert pool of an MoE transformer with \emph{deterministic computation
experts}, which are non-neural expert modules that execute exact domain
computations. Each deterministic expert comprises three stages:
\begin{enumerate}
\item A \textbf{parameter extraction MLP} (learned) that reads the hidden
state and outputs the numerical arguments required by the target
function.
\item A \textbf{deterministic computation} (exact) that executes the
domain formula as a standard PyTorch function, maintaining
differentiability through autograd.
\item A \textbf{result projection MLP} (learned) that maps the scalar
output back to the model's hidden dimension for re-injection into
the residual stream.
\end{enumerate}

We implement Hybrid MoE on top of OpenAI's gpt-oss-20b, a 21-billion
parameter MoE model with 32 experts per layer and top-4 routing. We
extend the router from 32 to 36 experts by adding 4 deterministic
experts for commercial real estate (CRE) underwriting metrics:
DSCR, LTV, Capitalization Rate, and Debt Yield. The entire modification
adds approximately 1.5 million parameters per extraction probe
(approximately 12 million total) while the base model remains fully
frozen.

Our contributions are:
\begin{enumerate}
\item \textbf{A novel heterogeneous expert architecture} that integrates
deterministic Python functions as first-class experts in MoE
routing, enabling exact computation within a single forward pass
(\Cref{sec:method}).
\item \textbf{Empirical demonstration} that financial numerical values
can be recovered from transformer hidden states with 5-fold
cross-validated $R^2$ up to 0.98, and that a learned router
achieves 100\% accuracy in detecting computation-requiring inputs
(\Cref{sec:experiments}).
\item \textbf{A 173$\times$ speedup} over generation-based methods
(40\,ms vs.\ 7{,}000\,ms per scenario), achieving 60--87\%
accuracy within 20\% error across all four experts with guaranteed
coverage ($N{=}200$), while generation baselines produce parseable
answers for only 1--94\% of scenarios (\Cref{sec:experiments}).
\item \textbf{Analysis of extraction robustness}, demonstrating that
probe training instability (not hidden state quality) causes
extraction failures on small datasets, and that 5-fold CV
ensembles resolve this (\Cref{sec:analysis}).
\end{enumerate}
